\chapter{Comparative Analysis}

\section{Chapter continuity}

This chapter continues the scientific positioning established in Chapter 3. Instead of introducing new theoretical concepts, it compares selected original research works according to criteria derived from the identified gaps: sensing, communication, prediction, simulation, recommendation, scheduling, actuation, validation, and deployment feasibility.

\section{Comparison criteria}

The comparison uses original research papers only. Review papers and overview papers are intentionally excluded from the matrix. The criteria are selected to reflect the functions needed in intelligent irrigation: sensing, communication, data integration, prediction, Digital Twin behavior, actuation, validation, and limitations.

\section{Functional comparison}

\begin{footnotesize}
\setlength{\tabcolsep}{2pt}
\renewcommand{\arraystretch}{1.15}
\begin{longtable}{@{}p{0.16\textwidth}p{0.18\textwidth}p{0.18\textwidth}p{0.18\textwidth}p{0.20\textwidth}@{}}
\caption{Functional comparison of selected original works}
\label{tab:functional-comparison-original-works}\\
\toprule
Study & Main focus & Strongest contribution & Digital Twin maturity & Limitation for this thesis \\
\midrule
\endfirsthead
\toprule
Study & Main focus & Strongest contribution & Digital Twin maturity & Limitation for this thesis \\
\midrule
\endhead
Kodali and Sarjerao \cite{kodali2017mqttirrigation} & Low-cost MQTT smart irrigation & Simple sensor-to-pump loop with remote status & Not a Digital Twin & Useful for MQTT/control background only \\
Balaceanu et al. \cite{balaceanu2019libelium} & Precision-agriculture monitoring & Multi-variable agricultural sensing & Monitoring layer only & No prediction, simulation, or actuation \\
Tancredi et al. \cite{tancredi2025irrigationdt} & ML-assisted irrigation decisions & Sensor-data classification for irrigation management & ML-enhanced Digital Twin direction & Limited field and implementation detail \\
Bellvert et al. \cite{bellvert2023vineyarddt} & Vineyard irrigation scheduling & Sentinel-2 assimilation and automated scheduling & Field-oriented irrigation Digital Twin & Requires remote-sensing and vineyard-specific calibration \\
Millan et al. \cite{millan2025tomatodt} & Tomato farm water-efficient management & Commercial-farm Digital Twin validation & Field-oriented irrigation Digital Twin & Site- and crop-specific management context \\
El Hachimi et al. \cite{elhachimi2025collectiveirrigationdt} & Collective irrigation-water distribution & Digital Twin with deep reinforcement learning & Optimization-oriented Digital Twin & Advanced optimization beyond prototype scope \\
Kallenberg et al. \cite{kallenberg2025interoperabledt} & Interoperable agricultural Digital Twins & FIWARE integration and RL recommendations & Interoperable AI-enhanced Digital Twin & General smart-farming scope, not irrigation-specific \\
Ed-Daoudi et al. \cite{eddaoudi2024predictiveirrigation} & Predictive smart irrigation & IoT and ML for adaptive irrigation scheduling & Predictive decision support & Limited explicit Digital Twin behavior \\
Gupta et al. \cite{gupta2025iotirrigation} & Automated IoT irrigation & Water and energy efficiency through automated control & IoT automation, not Digital Twin & No simulation or synchronized virtual state \\
Liu et al. \cite{liu2025cottondt} & Cotton water-stress Digital Twin & Crop-level water-stress representation & Crop Digital Twin & No direct irrigation-control workflow \\
\bottomrule
\end{longtable}
\end{footnotesize}

\section{Positioning of the thesis contribution}

The thesis contribution is positioned between low-cost IoT irrigation prototypes and field-validated Digital Twin irrigation systems. The low-cost prototypes provide useful principles for sensing, MQTT communication, and pump control \cite{kodali2017mqttirrigation}. Field Digital Twins provide the agronomic target: scheduling should be informed by data, models, and validation under real crop conditions \cite{bellvert2023vineyarddt,millan2025tomatodt}. Machine-learning and reinforcement-learning studies show how intelligence can be added, but they also confirm that prediction and optimization require careful validation \cite{tancredi2025irrigationdt,elhachimi2025collectiveirrigationdt,kallenberg2025interoperabledt}.

At this stage of the thesis, the contribution should be described only at a general level: an intelligent irrigation prototype that integrates monitoring, decision support, Digital Twin reasoning, and supervised control. Detailed architecture, technologies, database schema, model implementation, interfaces, and actuator behavior are reserved for the design and implementation chapters.

\section{Limitations to keep explicit}

The comparison indicates several limitations that must remain visible throughout the thesis. A prototype is not equivalent to a multi-season commercial field deployment. Prediction metrics do not automatically prove agronomic performance. Simulation must be calibrated before it can be used for operational water accounting. Physical control requires safety, authentication, and operator supervision. Finally, diagrams in the implementation chapters should be focused and modular so that each workflow can be understood and validated separately.

\section{Conclusion}

This chapter compared the selected original research papers and clarified the thesis position. The contribution is motivated by an integration gap rather than by a claim to outperform field-validated irrigation Digital Twins. The following chapters can now present requirements and design choices while keeping the scientific background cleanly separated from the implemented solution.
